\section{Solución Independiente}

\begin{flushleft}
\justifying
Tras analizar los problemas presentados en la Defensoría de los Derechos Politécnicos (DDP), se plantean diversas soluciones específicas para las problemáticas identificadas. 
Este análisis por cada solución asegura que cada módulo de forma independiente proporcione un valor agregado incluso antes de la integración total del sistema.

\subsection{Clasificador de Competencia}

Este componente funciona como un apoyo para manejar el alto volumen de quejas que recibe la Defensoría diariamente. Es importante destacar que este asistente no busca sustituir el criterio o el trabajo de los abogados, sino que actúa como una herramienta para aligerar su carga de trabajo. Al realizar un filtrado inicial, ayuda a identificar rápidamente si un caso le corresponde a la DDP o si debe enviarse a otra área. De esta manera se optimiza el proceso de recepción y se garantiza que los casos sean atendidos por el área correspondiente, reduciendo tiempos de espera y mejorando la eficiencia en la gestión de las quejas.

\subsection{Repositorio Digital de Información}

La implementación de un repositorio digital para almacenar la información de las quejas y los expedientes es una solución clave para abordar los problemas relacionados con la pérdida de archivos, el desgaste del papel y la duplicidad de folios. Este repositorio permitiría digitalizar toda la documentación, facilitando su acceso y gestión. Además de que al contar con una base de datos centralizada, se busca manejar los folios de forma automatizada para evitar las inconsistencias en los formatos de folio reduciendo así el riesgo de duplicidad de estos.

\subsection{Portal de seguimiento y comunicación}

Este módulo está dirigido al quejoso buscando resolver la problemática de la falta de comunicación con la DDP. Ya que el quejoso podría consultar el estatus de su queja gracias al sistema de folios digitales y el tablero de consulta, permitiéndole al usuario saber en qué parte del proceso se encuentra su queja. Esto reduce las dudas de los quejosos y mejora su experiencia al interactuar con la DDP, ya que no tendrían que depender de llamadas telefónicas o visitas presenciales para obtener información sobre el progreso de su caso.

\subsection{Mitigación de quejas incompletas}
Para resolver el problema de las solicitudes que llegan con datos faltantes, se plantea un formulario con validaciones. Gracias a las reglas de validación obligatorias, se garantiza que cada expediente cuente con la información y los documentos requeridos antes de ser procesado. Esto permite que los abogados asesores reciban casos bien estructurados desde el inicio, eliminando la necesidad de contactar repetidamente al quejoso para subsanar omisiones, lo que agiliza drásticamente el análisis jurídico.

\subsection{Control de Acceso y Privacidad de la Información}

Debido a que el personal puede cambiar constantemente, este componente basado en un sistema de roles, asegura que la información sensible esté siempre protegida. Su función es organizar quién puede ver o modificar cada parte de la queja según su rol en la DDP. Esto ayuda a que los nuevos integrantes del equipo aprendan a usar la plataforma más rápido, ya que solo interactúan con las herramientas que necesitan para su trabajo diario, manteniendo siempre un control claro sobre la privacidad de los datos de la comunidad.

Cada una de estas soluciones de forma independiente aporta beneficios específicos tanto para la DDP como para los quejosos, mejorando la eficiencia, la organización y la comunicación dentro de la institución. 
Sin embargo, es importante mencionar que al integrar estas soluciones en una sola plataforma unificada potencia los beneficios dados, como se observará en el siguiente capítulo.
\end{flushleft}